% Окружения утверждений (ручной файл; основа -- hse-uberhandbook/common/theorems.tex).
% Отличия от uberhandbook: нумерация по главам («теорема 1.2», как в конспекте);
% окружения proposition («Утверждение») и method («Метод»); example и remark без
% бокса; список утверждений (\tcblistof{thms}); \proofpart / \proofstep для шагов
% доказательства; пустой заголовок не оставляет висячего двоеточия.
\usepackage{amsthm}

\usepackage[most]{tcolorbox}
\tcbuselibrary{breakable,skins}

\usepackage{xparse}

\newcounter{theorem}[chapter]
\renewcommand{\thetheorem}{\thechapter.\arabic{theorem}}

\tcbset{
  thmbox/.style={
    before skip=8pt,
    after skip=10pt,
    enhanced,
    breakable,
    colback=white,
    colframe=black,
    boxrule=0.8pt,
    arc=0pt,
    left=2mm, right=2mm, top=1.5mm, bottom=1.5mm,
    boxsep=1.2mm,
    title filled,
    colbacktitle=gray!15,
    coltitle=black,
    fonttitle=\normalfont,
  },
  methodbox/.style={
    thmbox,
    colbacktitle=blue!8,
  },
}

% #1 -- вид («Теорема»), #2 -- заголовок, #3 -- стиль бокса
\NewDocumentEnvironment{_boxedthm}{ m O{} O{thmbox} }
{%
  \refstepcounter{theorem}%
  \addcontentsline{thms}{tcolorbox}{\protect\numberline{\thetheorem}#1\IfBlankF{#2}{: #2}}%
  \begin{tcolorbox}[
    #3,
    title={%
      {\bfseries #1~\thetheorem}%
      \IfBlankF{#2}{:\ {\normalfont\itshape #2}}%
    },
  ]%
}
{%
  \end{tcolorbox}%
}

\NewDocumentEnvironment{theorem}{ O{} }{\begin{_boxedthm}{Теорема}[#1]}{\end{_boxedthm}}
\NewDocumentEnvironment{lemma}{ O{} }{\begin{_boxedthm}{Лемма}[#1]}{\end{_boxedthm}}
\NewDocumentEnvironment{corollary}{ O{} }{\begin{_boxedthm}{Следствие}[#1]}{\end{_boxedthm}}
\NewDocumentEnvironment{definition}{ O{} }{\begin{_boxedthm}{Определение}[#1]}{\end{_boxedthm}}
\NewDocumentEnvironment{proposition}{ O{} }{\begin{_boxedthm}{Утверждение}[#1]}{\end{_boxedthm}}
\NewDocumentEnvironment{method}{ O{} }{\begin{_boxedthm}{Метод}[#1][methodbox]}{\end{_boxedthm}}

% Небоксовые: пример (нумеруется, попадает в список) и замечание (без номера).
% #1 -- вид, #2 -- заголовок
\NewDocumentEnvironment{_plainthm}{ m O{} }
{%
  \refstepcounter{theorem}%
  \addcontentsline{thms}{tcolorbox}{\protect\numberline{\thetheorem}#1\IfBlankF{#2}{: #2}}%
  \par\addvspace{8pt}%
  \noindent\textbf{#1~\thetheorem}%
  \IfBlankF{#2}{:\ {\normalfont\itshape #2}}%
  \ \normalfont\ignorespaces
}
{%
  \par\addvspace{10pt}%
}
\NewDocumentEnvironment{example}{ O{} }{\begin{_plainthm}{Пример}[#1]}{\end{_plainthm}}
\NewDocumentEnvironment{concept}{ O{} }{\begin{_plainthm}{Понятие}[#1]}{\end{_plainthm}}
\NewDocumentEnvironment{remark}{ O{} }
{%
  \par\addvspace{6pt}%
  \noindent\textbf{Замечание}%
  \IfBlankF{#1}{\ {\normalfont\itshape (#1)}}%
  .\ \normalfont\ignorespaces
}
{%
  \par\addvspace{8pt}%
}
\NewDocumentEnvironment{boxedexample}{ O{} }{\begin{_boxedthm}{Пример}[#1]}{\end{_boxedthm}}

\newtheorem*{proofsketch}{Набросок доказательства}

% Шаги доказательства: \proofpart{⇒} -- крупный шаг, \proofstep{…} -- подшаг.
\NewDocumentCommand{\proofpart}{m}{\par\medskip\noindent{\itshape #1.}\par\nobreak\smallskip\noindent\ignorespaces}
\NewDocumentCommand{\proofstep}{m}{\par\smallskip\noindent{\itshape #1.}\enspace\ignorespaces}

% Повтор заметки в другой главе: ссылка на первый экземпляр вместо второго.
% #1 -- id заметки, #2 -- заголовок
\NewDocumentCommand{\noteagain}{m m}{%
  \par\addvspace{6pt}\noindent
  {\itshape См.}\ \hyperref[note:#1]{\textbf{#2}} (стр.~\pageref{note:#1}).%
  \par\addvspace{8pt}}
